П'ять основ — це фундамент учення. Їх не можна ігнорувати, спотворювати чи розглядати без належної уваги. Насправді п'ять основ — це \emph{цілком істотні складники} вчення, які застосовуються як до пробудження, так і, можливо, ще сильніше, після нього. Не обманюйте себе думкою, що п'ять основ є незначними чи примітивними лише тому, що вони, здається, узгоджуються тільки з людським, або відносним, аспектом Дійсності. П'ять основ — це засіб \emph{жити} й \emph{виявляти} останню природу Дійсності в повсякденні. Якщо ми не живемо й не виявляємо у своєму житті те, що ми здійснили в найглибші моменти одкровення, то ми житимемо розділеним життям.

Крім того, п'ять основ задають контекст, у якому розгортається вчення. Якщо ви вилучите контекст учення, ви вилучите й антиегоїстичні запобіжники, які захищають учення від самозосереджених тлумачень. Хибне тлумачення духовного вчення его завжди є значною небезпекою, адже схильність его — виправдовувати будь-яку точку зору, у яку воно вклалося і до якої прив'язалося.

На додаток до цієї небезпеки, будь-яке духовне вчення, вкорінене в абсолютній природі Дійсності, за визначенням зорієнтоване на Істину, а не на етичні й моральні виміри життя. Це не означає, що такі вчення аморальні; вони трансморальні, тобто вкорінені в дійсності, що перебуває за межами моралі й етики, які стосуються взірців дуалістичної перспективи.

Не йдеться про те, що мораль не має значення в абсолютному баченні; це досить поширене хибне тлумачення. Ідеться про те, що мораль більше не ґрунтується на культурних і релігійних цінностях, призначених стримувати егоїстичні імпульси. Натомість безкорислива любов і співчуття природно випливають з об'єднаного бачення Дійсності як спонтанного вияву Єдності. Коли ніщо не вбачається як відокремлене чи як щось, що не є вами, дії, які випливають із вас, є відображенням об'єднаної перспективи.

Можливі певні ускладнення, адже можна мати \emph{деякий} досвід останньої Дійсності, коли ще не сталося повного звільнення від егоїстичної омани. Це відбувається через можливу летку суміш Дійсності й ілюзії, які одночасно виявляються й існують у несвідомий і викривлений спосіб. Хоча деякі такі випадки можна очікувати в міру духовного дозрівання, є речі більш викривлені й небезпечні, ніж его, яке вважає себе Богом.

Багато років роботи з тисячами людей показали мені, що якщо цими засадничими аспектами духовного життя знехтувати, духовні одкровення дуже ймовірно зійдуть з рейок у той чи інший спосіб. Невдача в дослідженні й чіткому розумінні будь-якої з них, а також у послідовному їх застосуванні призведе до якоїсь форми конфлікту, поділу й сум'яття.

П'ять основ — це засоби зібрати всі свої внутрішні ресурси — тіло, розум і дух — і зосередити їх об'єднано задля свого найвищого прагнення. Я не можу достатньо наголосити на важливості ясного й об'єднаного фокусу, щирого серця й непохитного прагнення не обманювати ні себе, ні інших.

\section*{Проясніть своє прагнення}

Прояснити своє прагнення означає точно знати мотивацію свого духовного життя, \emph{не як мету на майбутнє, а на кожну мить}. Іншими словами — те, що ви найбільше цінуєте в житті, не в сенсі моральних цінностей, а в сенсі того, що для вас найважливіше. Обміркуйте це питання. Не припускайте просто, що ви вже знаєте своє найвище прагнення чи те, що для вас найважливіше. Заглибтеся й подумайте над тим, що найважливіше у вашому духовному пошуку; не дозволяйте нікому визначити це за вас. Не заспокоюйтеся, доки вам не стане цілком ясно, яке ваше прагнення.

Важливість цієї першої основи слід добре наголосити, адже життя розгортається на основі того, що ви найбільше цінуєте. Дуже мало людей мають Істину або Дійсність як глибинні цінності. Багато хто думає, що цінує Істину, але їхні дії цьому не відповідають. Загалом люди мають конкуруючі й суперечливі цінності, що виявляється у внутрішніх і зовнішніх конфліктах. Отже, те, що ви \emph{думаєте}, ніби щось є вашою найглибшою цінністю, не означає, що воно справді нею є. Обмірковуючи й прояснюючи те, що ви цінуєте й до чого прагнете, ви стаєте більш об'єднаним, ясним і впевненим у своєму напрямку.

У міру того, як ваше здійснення й духовна зрілість поглиблюються, ви помітите, що деякі аспекти вашого життя залишаються сталими, тоді як інші розвиваються, щоб відображати те, що є доречним для вашого рівня розрізнення. Так створюються більший фокус і зв'язок із перебігом вашого власного здійснення.

\section*{Безумовна послідовність}

Прояснення свого прагнення — це перший крок. Так енергія й увага зводяться в об'єднану силу, яка рухає вас до вашого прагнення. Коли ваше прагнення прояснене, ви повинні рухатися до нього. Така постава стосується того, що ви \emph{готові робити}, і передусім того, від чого ви \emph{готові відмовитися}.

Духовність не вимагає, щоб ви тяжко працювали задля якогось майбутнього результату — вона вимагає, щоб ви були цілком присутні, щирі й віддані \emph{тепер}, з абсолютною чесністю й готовністю виявити й відпустити \emph{будь-яку} ілюзію, що стоїть між вами та здійсненням Істини. Отже, духовність не має нічого спільного з часом чи з тим, чого можна досягти з часом. Духовність стосується \emph{лише й завжди} вічного теперішнього.

\emph{Прагнення} — це радше питання серця, ніж розуму, адже йдеться про те, що ви найбільше цінуєте, любите й бережете. Вам не треба нагадувати про те, що ви любите, а про те, чого ви \emph{не} любите. І те, що ви \emph{справді} любите, найправдивіше відображається у ваших \emph{діях}, а не в тому, що ви відчуваєте, думаєте чи говорите.

Коли прагнення узгоджується з безумовною послідовністю та любов'ю, воно стає вельми потужною силою у всесвіті. Лише так ми будемо достатньо об'єднаними й зосередженими, щоб наше прагнення пережило можливі обставинні труднощі.

\section*{Ніколи не зрікайтеся свого авторитету}

Третя основа — \emph{ніколи не зрікатися свого авторитету}. Це означає, що \emph{ви будете цілком відповідальні за своє життя і ніколи не поступитеся цим жодній людині}. Що ж до самого процесу просвітлення, то безглуздо хотіти прокататися на спині іншої просвітленої істоти. Будь-яка невдача в розумінні цієї основи може призвести (як і призводила в багатьох випадках) до культу фанатизму, фундаменталізму, містичного мислення, до розчарування, омани та/або духовної незрілості.

Зрозуміло, що багато людей проєктують свої сімейні, стосункові, сексуальні проблеми, а також проблеми з Богом на своїх духовних учителів (і часто несумлінні вчителі заохочують їх до цього), але істотно зрозуміти, що роль духовного вчителя — бути добрим і мудрим духовним провідником, утіленням Істини, на яку він чи вона вказує. Може бути навіть глибока повага, любов або й відданість духовному вчителеві, але засадничо важливо не зрікатися власного авторитету на користь його авторитету й не проєктувати всю божественність виключно на нього. Ваше життя належить вашим рукам, а не чиїмось іншим. Будьте належно відповідальними.

Існує дуже тонка межа між тим, щоб бути справді відкритим до настанов духовного вчителя, і регресом до дитячих стосунків, у яких ви зрікаєтеся свого авторитету й проєктуєте \emph{всю} мудрість і божественність на вчителя. Усім людям потрібно знайти зрілу рівновагу у відкритості до духовних настанов, не зрікаючись власного авторитету.

Те саме можна застосувати й до духовного вчення. Кожне духовне вчення — це палець, що вказує на Дійсність, але ніколи не сама Дійсність. Перебування у зрілих стосунках із духовним ученням вимагає, щоб ви його \emph{застосовували}, а не просто вірили в нього. Вірування ведуть до різних форм фундаменталізму й заглушують цікавість та дослідження, які є істотними, щоб дати простір пробудженню і тому, що є за межами пробудження. Добре духовне вчення — це те, що втілюють у \emph{практику}, це те, з чим ви \emph{працюєте}. Роблячи це, воно працює у вас (зазвичай потаємно) і допомагає вам виявити Істину (і неправди), що дрімають у вас.

Що ж означало б не зрікатися свого авторитету й при цьому не проголошувати хибний чи самозосереджений авторитет, який приведе вас до омани? Боюся, що не можу цього сказати. Ніхто не може сказати вам, як не обманювати себе. Якщо в найглибшому у вас є прагнення Істини понад усе, то навіть якщо ви заблукаєте тисячами способів, ви завжди знову знайдете себе — вас буде повернено до того, що є Істинним.

А якщо ви не прагнете Істини понад усе — що ж, ви вже знаєте, куди це вас приведе.

\section*{Практикуйте абсолютну щирість}

Мати справжню щирість — абсолютно необхідно в духовному житті. Щирість охоплює такі якості, як чесність, справжність і цілісність. Бути щирим не означає бути досконалим. Насправді саме зусилля бути досконалим не є щирим, бо це ще один спосіб уникнути бачення себе таким, яким ви є саме тепер. Здатність і бажання бачити себе таким, яким ви насправді є, з усіма своїми недосконалостями й ілюзіями, вимагає справжньої щирості й мужності. Якщо ми постійно ховаємося від себе, ми ніколи не зможемо прокинутися від ілюзії его.

Щоб бути щирими, нам потрібно перестати засуджувати себе. Осуд затуляє доступ до справжньої щирості й навіть може маскуватися \emph{під} щирість. Справжня щирість виявляє потужну форму ясності й розрізнення, яка потрібна, щоб ми могли чесно бачити себе без утечі й без ув'язнення в зумовлених судженнях розуму та його захисних механізмах.

Здатність і бажання бути чесним із собою — ваш найбільший щит проти самосаботажу, і воно узгоджує вас із вашим справжнім прагненням. Немає більшого виклику для людини, ніж бути цілком чесною із собою, а також з іншими. Така чесність абсолютно необхідна, якщо ми справді прагнемо прокинутися від сну відокремленості й жити справжнім та об'єднаним життям.

\section*{Добре керуйте своїм життям}

Добре керувати власним життям означає не використовувати духовність, щоб уникати якогось аспекту свого життя. Я спостерігав, що серед людей, залучених до духовності, дуже поширене несвідоме використання її для уникнення болісних, проблемних, дисфункційних або страшливих аспектів себе. Часто є надія, що якщо вони просто прокинеться до Дійсності, усі виклики зникнуть. Хоча правда, що з появою пробудження багато аспектів, які ми вважаємо проблемами, просто зникнуть, було б помилкою припускати, що проблиск неспанного стану автоматично розв'яже всі складні аспекти людського життя.

Використання духовності для уникнення складних аспектів себе чи свого повсякдення може значною мірою стримати появу духовного просвітлення і неодмінно стримає його глибину й стабільність. Шлях до звільнення — це спосіб подивитися на себе й своє життя \emph{цілком}, не втікаючи в заперечення, в осуд чи в магічне мислення. Це спосіб розрізати завіси ілюзії й прокинутися до Істини.

Бути добрим управителем життя вимагає, щоб ви прийняли кожен аспект свого життя — внутрішній і зовнішній, приємний чи неприємний. Вам не обов'язково зустрічати їх усі одразу, лише те, що постає в цю мить. Дайте кожній миті увагу, щирість і відданість, на які вона заслуговує. Невдача в цьому коштує набагато дорожче, ніж ви можете собі уявити.

Ваше життя, \emph{все} ваше життя — це ваш шлях до пробудження. Чинячи опір своїм викликам або не маючи з ними справи, ви триматимете себе в заціпенінні щодо Дійсності. Звертайте увагу на те, що життя намагається вам відкрити. Скажіть «так» його силі, його суворості та його люблячій благодаті.
